\runninghead{Scope and References}
\section{Appendix: Scope and Qualifications}
\label{sec:scope}

{\small
\textbf{Edition, Mass and occurrence.} The study treats the Mass
\latin{Dominica prima Adventus}, I classis, \latin{Statio ad S. Mariam
maiorem}, of the \work{Missale Romanum}, editio typica, Typis Polyglottis
Vaticanis 1962, printed pp.~1--2, nos.~1--10, as said on Sunday 29~November
2026 in the universal 1962 calendar and no particular calendar. The First
Sunday of Advent is the Sunday falling on or nearest to 30~November (RG~20),
and it is of the first class (RG~11~a). The commemoration of St Saturninus,
which the calendar sets on 29~November, is an ordinary commemoration and is
not admitted on a day of the first class (RG~107--109, 111~a; RGMR~434~a).
The colour is violet (RG~128~a); the Gloria is not said (the rubric before
the Collect and RGMR~432~b); the Credo is said (the rubric after no.~6 and RGMR~475~a); the
Preface is of the Most Holy Trinity (the direction after no.~8 and
RGMR~494~b), as \texttt{research/context.md} records. Two dated Ordos of
institutes using the 1962 books, of the Fraternity of St Peter and of the
Institute of Christ the King Sovereign Priest, both in France, give the same
Mass for the date, of the first class and in violet. No local oration,
external solemnity or votive Mass is considered. The Ordinary and the Office
are not treated.

\textbf{Text control.} Every Latin form was read on page images of the CMAA
facsimile of the 1962 typical edition and collated word for word with the
Ratisbon printing of 1862 (Pustet), pp.~1--3, whose page images carry every
word of the ten texts; the collation and the comparison with the Clementine
Vulgate are in \texttt{propers/verified.md}. Where the chants and lessons
differ from the Clementine Vulgate the differences are reported as textual
facts. Augustine's psalter agrees with the chants in \latin{qui te
exspectant} and \latin{notas fac}, and differs from them in \latin{Ad te,
Domine, levavi} and \latin{dabit suavitatem}; the Latin psalter or
antiphonary from which the chant forms came was not identified. The history
of these texts' assignment to this Sunday is not treated.

\textbf{English.} The English of the lessons and chants is the Douay--Rheims
in Challoner's revision, with the Missal's departures from the Vulgate named
beneath each text; the English of the three orations and of the Communion is
that of the Lasance \work{New Roman Missal} in the Benziger revision of 1945,
public domain in the United States by non-renewal. None of this English is
an approved liturgical translation of the 1962 Missal. Chrysostom on Romans
and on Matthew, Augustine on Psalm 24 and the \work{Confessions}, and Cyril of
Jerusalem are quoted in the English of the Nicene and Post-Nicene Fathers;
Cyril of Alexandria in R.~Payne Smith's English of 1859; Bellarmine in John
O'Sullivan's English of 1866, which its translator abridged; Schuster in the
English of 1924. Elsewhere, and wherever a Father's own Latin is quoted, the
English beside it is a rendering made for this study. Where the argument
turns on the Latin of a chant or an oration, the commentary renders the words
or the clause in question literally within its own sentences; these
renderings are not offered as a version of the prayer, whose English is the
Lasance or Douay text printed with it.

\textbf{Reception.} The interpretations rest on the passages named in the text
and recorded, with editions and bounds, in \texttt{research/scope.md} and
\texttt{research/interpretations.md}. Every connection drawn between two texts
of the Mass belongs to the study's interpretations unless an author is named as
making it; no Father cited here comments on this Mass as a whole. Gregory the
Great, Ambrose, Cyril of Alexandria, Augustine's Letter~199 and Bede comment on
Luke's discourse; Chrysostom on Matthew and Cyril of Jerusalem's
\work{Catechetical Lectures} comment on the parallel passages. Gregory's text
of Lk~21:27 reads \latin{in nubibus}, where the Missal has \latin{in nube}.
Augustine's \work{Confessions} narrate his
conversion and are not an exposition of Romans. Gregory, Augustine's
\work{Enarrationes} on Psalm 84 and Ambrose are read in Migne's Latin as
transcribed on Wikisource and in the Corpus Corporum, not collated with page
images; Augustine's Letter~199 and his \work{Enarratio} on Psalm~24 in the
augustinus.it text; the Latin of the \work{Confessions} in J.~J. O'Donnell's
text as given by The Latin Library; Aquinas on Romans and
on Psalm 24 in the Corpus Thomisticum text; Rupert, Sicard and Durandus in
uncorrected optical text, whose quotations here are normalised in spelling.
Every principal author of an interpretation has a row in the repository's
register of authorial standing. Rupert of Deutz, Sicard of Cremona and William
Durandus are ecclesiastical writers and Blessed Ildefonso Schuster is
Blessed; none of them is a principal author. Rupert and Sicard expound a Mass
of this Sunday whose Gospel was another, and they are cited for the chants
and the Epistle it shared with this one; Sicard is also cited for what he adds
of the churches that read \latin{Erunt signa}, which Durandus copies. Cassiodorus on the psalms, Jerome on
Matthew, Augustine's \work{City of God}~XX and the \work{Catena aurea} were
not read.

\textbf{Chronology.} Every biblical date in the Scriptural Date and Location
appendix is generated from the repository's Scripture chronology corpus,
with its sources named there.

\textbf{Rights.} The Latin is published on 17 U.S.C. 103(b) together with its
public-domain witness, the Pustet printing of 1862. The Douay--Rheims and the
Lasance English are quoted from registered editions. The other quotations are
short and come from editions published before 1931 or from ancient and
medieval Latin: Migne's \work{Patrologia Latina}, vols.~15, 92, 170 and 213,
and the Migne texts of Augustine and Gregory; the Lyon \work{Rationale} of
1612; the Nicene and Post-Nicene Fathers; Payne Smith's Cyril (1859);
O'Sullivan's Bellarmine (1866); and Schuster's \work{Sacramentary} (1924). The
Corpus Thomisticum, augustinus.it and Wikisource texts are used only for Latin
that is itself medieval or ancient. Scripture, liturgical texts and quoted
translations are not relicensed by the repository's licence.

\textbf{Use.} This is a study aid to the 1962 Missal, not an official
liturgical book, translation or critical edition.
\par}

\section{References}
\label{sec:references}

{\small
\subsection*{Liturgical books}
\begin{itemize}[leftmargin=1.2em,itemsep=.25em]
\item \work{Missale Romanum ex decreto Sacrosancti Concilii Tridentini
restitutum, Summorum Pontificum cura recognitum}, editio typica (Typis
Polyglottis Vaticanis, 1962), CMAA facsimile: \work{Rubricae generales} and
\work{Rubricae generales Missalis romani}; Proprium de Tempore, pp.~1--2,
nos.~1--10.
\item \work{Missale Romanum} (Ratisbon: Pustet, 1862), pp.~1--3.
\item Lasance, F.~X., and Walsh, F.~A., \work{The New Roman Missal in Latin and
English}, revised edition (New York: Benziger Brothers, 1945), ``First Sunday
of Advent,'' pp.~103--106.
\item Fraternité Saint-Pierre (France), \work{Ordo du mois}, and Institut du
Christ Roi Souverain Prêtre (France), \work{Ordo}, entries for 29~November
2026 (websites), read 25~November 2026.
\end{itemize}

\subsection*{Scripture}
\begin{itemize}[leftmargin=1.2em,itemsep=.25em]
\item \work{The Holy Bible Translated from the Latin Vulgate} (Douay--Rheims,
Challoner revision), Project Gutenberg text.
\item \work{Biblia Sacra Vulgatae Editionis} (Clementine Vulgate), for the
comparisons with the Missal's forms.
\end{itemize}

\subsection*{Fathers, Doctors and other authors}
\begin{itemize}[leftmargin=1.2em,itemsep=.25em]
\item Ambrose, \work{Expositio evangelii secundum Lucam} X.37--42, Migne
Latin text, PL~15 (Wikisource).
\item Augustine, \work{Enarrationes in Psalmos} 24.1--4 and 24.10, NPNF,
first series, vol.~8; 24.2--4, augustinus.it Latin text; 84.9 and 84.15,
Migne Latin text (Wikisource); \work{Epistula} 199 to Hesychius,
augustinus.it Latin text; \work{Confessions} VIII.5.12 and VIII.12.29,
trans.\ J.~G. Pilkington, NPNF, first series, vol.~1, and VIII.5.12, Latin
text ed.\ J.~J. O'Donnell (The Latin Library).
\item Bede, \work{In Lucae evangelium expositio} VI, on 21:32--33, Migne Latin
text (Corpus Corporum).
\item Bellarmine, Robert, \work{Explanatio in Psalmos}, on Ps~24:2--4 and
Ps~84:8 and 13, trans.\ John O'Sullivan, \work{A Commentary on the Book of
Psalms} (Dublin: Duffy, 1866).
\item Chrysostom, John, \work{Homilies on Romans} 24, NPNF, first series,
vol.~11; \work{Homilies on Matthew} 77, NPNF, first series, vol.~10.
\item Cyril of Alexandria, \work{Commentary on the Gospel of Saint Luke},
sermon 139, trans.\ R.~Payne Smith (Oxford, 1859).
\item Cyril of Jerusalem, \work{Catechetical Lectures} 15.1, 15.3--4, NPNF,
second series, vol.~7.
\item Gregory the Great, \work{Homiliae in Evangelia} I.1, nn.~1--6, Migne
Latin text (Wikisource).
\item Thomas Aquinas, \work{Super Epistolam ad Romanos}, cap.~13, lect.~3,
and \work{Super Psalmos}, on Ps~24, nn.~1 and 4, Corpus Thomisticum text.
\end{itemize}

\subsection*{Commentators on the Mass}
\begin{itemize}[leftmargin=1.2em,itemsep=.25em]
\item Durandus, William, \work{Rationale divinorum officiorum} VI, on the
first Sunday of Advent (Lyon, 1612).
\item Rupert of Deutz, \work{De divinis officiis} III.2, PL~170.
\item Schuster, Ildefonso, \work{The Sacramentary (Liber Sacramentorum)},
vol.~1 (London: Burns, Oates and Washbourne, 1924), ``First Sunday of
Advent,'' pp.~321--323.
\item Sicard of Cremona, \work{Mitrale} V, on the first Sunday of Advent,
PL~213.
\end{itemize}
\par}
